\clearpage
\begin{examplebox}
\textbf{例1：独立核验首页图中的两条平行线}

已知 $C:x^2/2-y^2=1$、$P(2,1)$，两割线斜率为 $2,-2$。求第二交点并验证结论。

\textbf{先算交点：} 设 $x=2+h$。第一条直线 $y=1+2h$ 代入得
\[
h\left(-2-\frac72h\right)=0.
\]
舍去 $h=0$（旧点 $P$），得 $h=-4/7$，所以 $M=(10/7,-1/7)$。第二条直线 $y=1-2h$ 同理给出
\[
h\left(6-\frac72h\right)=0,
\]
第二个根为 $12/7$，所以 $N=(26/7,-17/7)$。

\textbf{再查方向和位置：}
\[
x_M+y_M=x_N+y_N=\frac97,
\qquad MN:\ x+y=\frac97.
\]
镜像点 $Q=(2,-1)$ 的切线为 $x+y=1$。两条直线斜率都为 $-1$，截距不同，确实平行。用结论只求方向时，就能省去前面求两点的计算。
\end{examplebox}

\clearpage
\begin{examplebox}
\textbf{例2：内心在竖直线上，截距为什么是开区间？}

$C:x^2/2-y^2=1$，$P(2,1)$，直线 $l$ 与 $C$ 交于不同的 $M,N$。若 $\triangle PMN$ 的内心在 $x=2$ 上，求 $l$ 在 $y$ 轴上的截距范围。

\textbf{识别：} 内心与顶点同在 $x=2$ 上，故两边射线关于它对称，割线斜率为 $m,-m$。由 C095，$l$ 平行于 $Q(2,-1)$ 点切线，因此设 $l:x+y=t$。

联立得交点横坐标方程 $f(x)=x^2-4tx+2t^2+2=0$。
\begin{center}
\begin{tikzpicture}[x=1cm,y=1cm,font=\small]
\draw[gray!45,rounded corners=3pt,fill=white] (0,0) rectangle (13.8,8.7);
\node[anchor=west,font=\small\bfseries] at (.4,8.2) {内心横坐标为 2，必须位于两交点横坐标之间};

\node[anchor=west] at (.6,7.55) {$t=1$：$f(x)=(x-2)^2$，两个根重合};
\draw[cnaxis] (1,6.7)--(12.4,6.7);
\fill[orange!85!black] (3,6.7) circle(2pt) node[below=5pt] {$2$（重根）};
\node[anchor=west] at (.6,5.65) {$t=2$：两根严格夹住 $2$，所以 $f(2)<0$};
\draw[cnaxis] (1,4.7)--(12.4,4.7);
\draw[green!45!black,very thick] ({5-sqrt(6)},4.7)--({5+sqrt(6)},4.7);
\fill ({5-sqrt(6)},4.7) circle(2pt) node[above left=4pt] {$4-\sqrt6$};
\fill ({5+sqrt(6)},4.7) circle(2pt) node[above=5pt] {$4+\sqrt6$};
\fill[red!78!black] (3,4.7) circle(2pt) node[below=5pt] {$2$};
\node[anchor=west] at (.6,3.55) {$t=3$：$f(x)=(x-2)(x-10)$，一个交点回到 $P$};
\draw[cnaxis] (1,2.6)--(12.4,2.6);
\draw[green!45!black,very thick] (3,2.6)--(11,2.6);
\fill[orange!85!black] (3,2.6) circle(2pt) node[below=5pt] {$2$};
\fill (11,2.6) circle(2pt) node[below=5pt] {$10$};
\node at (6.9,.8) {$f(2)=2(t-1)(t-3)<0\quad\Longleftrightarrow\quad 1<t<3$};

\end{tikzpicture}

\end{center}

内心在三角形内部，其横坐标 $2$ 必须严格夹在两根之间，所以
\[
f(2)=2(t-1)(t-3)<0\quad\Rightarrow\quad 1<t<3.
\]
\end{examplebox}

\clearpage
\begin{examplebox}
\textbf{例2（续）：区间内的每个值都能取到吗？}

当 $1<t<3$ 时，$f(2)<0$，二次函数开口向上，故有两个不同实根 $x_M<2<x_N$。由韦达定理，
\[
x_M+x_N=4t,\qquad
(x_M-2)(x_N-2)=2(t-1)(t-3).
\]
两条边的斜率分别为
\[
k_{PM}=\frac{t-x_M-1}{x_M-2}=\frac{t-3}{x_M-2}-1,
\quad
k_{PN}=\frac{t-3}{x_N-2}-1.
\]
因此
\[
\begin{aligned}
k_{PM}+k_{PN}
&=(t-3)\frac{x_M+x_N-4}{(x_M-2)(x_N-2)}-2\\
&=(t-3)\frac{4(t-1)}{2(t-1)(t-3)}-2=0.
\end{aligned}
\]
两根位于 $2$ 两侧，故这两条边射线关于竖直轴互为镜像，$x=2$ 是内角平分线。三点不共线（$t\ne3$），内心确实在 $x=2$ 上，所以所有 $1<t<3$ 均可取到。

\medskip
\textbf{答案与端点回验：} 截距范围为 $\boxed{(1,3)}$。$t=1$ 时两交点合并为 $Q$；$t=3$ 时一个交点回到 $P$，均不能形成题设三角形。图中的两端都必须排除。
\end{examplebox}

\clearpage
\begin{examplebox}
\textbf{例3：顶点处不用有限斜率，照样判断平行}

$C:x^2/2-y^2=1$、$P(\sqrt2,0)$，两割线斜率为 $2,-2$。求弦方向。

\begin{center}
\begin{tikzpicture}[x=1cm,y=1cm,font=\small]
\draw[gray!45,rounded corners=3pt,fill=white] (0,0) rectangle (13.8,8.2);
\node[anchor=west,font=\small\bfseries] at (.4,7.699999999999999) {顶点情形：方向存在，但不能再用有限斜率表示};

\begin{scope}[shift={(1.6,3.9)},scale=2]
\draw[cncurve,samples=120,domain=-1.2:1.2,variable=\t] plot ({sqrt(2*(1+\t*\t))},\t);
\draw[cntangent] ({sqrt(2)},-1.3)--({sqrt(2)},1.3);
\draw[cnchord] ({9*sqrt(2)/7},-1.25)--({9*sqrt(2)/7},1.25);
\draw[cnsec] ({9*sqrt(2)/7},{4*sqrt(2)/7})--({sqrt(2)},0)--({9*sqrt(2)/7},{-4*sqrt(2)/7});
\fill ({sqrt(2)},0) circle(1.8pt) node[left=6pt] {$P=Q$};
\fill ({9*sqrt(2)/7},{4*sqrt(2)/7}) circle(1.8pt) node[right=5pt] {$M$};
\fill ({9*sqrt(2)/7},{-4*sqrt(2)/7}) circle(1.8pt) node[right=5pt] {$N$};
\end{scope}
\node[anchor=west] at (8.4,6.3) {$P=Q=(\sqrt2,0)$};
\node[anchor=west,orange!85!black] at (8.4,4.9) {$\ell_Q:\ x=\sqrt2$};
\node[anchor=west,red!78!black] at (8.4,3.4) {$MN:\ x=\dfrac{9\sqrt2}{7}$};
\node[anchor=west] at (8.4,1.9) {都竖直，仍然平行};
\node at (6.9,.65) {本图仍取两条割线斜率为 $2,-2$；不写不存在的斜率分式。};

\end{tikzpicture}

\end{center}

\textbf{直接判断：} $Q=P$，此处切线为 $x=\sqrt2$，所以 $MN$ 竖直。

\textbf{独立回验：} 令 $x=\sqrt2+h$、$y=\pm2h$，代入并舍去 $h=0$，得 $h=2\sqrt2/7$。因此
\[
M=\left(\frac{9\sqrt2}{7},\frac{4\sqrt2}{7}\right),\quad
N=\left(\frac{9\sqrt2}{7},-\frac{4\sqrt2}{7}\right).
\]
两点横坐标相同，弦确实竖直；不能把不存在的有限斜率写成“无穷大”。
\end{examplebox}
